\documentclass{article}
\usepackage[utf8]{inputenc}
\usepackage{amsmath}
\usepackage{graphicx}
\usepackage{booktabs}
\usepackage{hyperref}

\title{Neural Stochastic Differential Equations for Option Pricing: \\
A Superior Alternative to Black-Scholes}
\author{NeuralBlackScholes Project}
\date{February 2026}

\begin{document}

\maketitle

\begin{abstract}
The Black-Scholes-Merton (BSM) model has long been the standard for option pricing, despite its restrictive assumption of constant volatility. In this paper, we present a Neural Stochastic Differential Equation (Neural SDE) model that learns underlying asset dynamics directly from market data. Trained utilizing the Adjoint Sensitivity Method for memory efficiency, our model demonstrates significant improvement in pricing accuracy. Benchmark results on S\&P 500 (SPY) options show a Mean Squared Error (MSE) of 0.26 compared to 389.87 for the BSM baseline, effectively capturing the volatility smile and term structure that analytical models often fail to address.
\end{abstract}

\section{Introduction}
Financial markets exhibit complex behaviors such as stochastic volatility, jumps, and heavy tails, which contradict the geometric Brownian motion assumption of the Black-Scholes model. While practitioners often employ implied volatility surfaces to compensate for these deficiencies, this approach is reactive rather than predictive. 

We propose a generative approach using Neural SDEs, which combine the mathematical rigor of differential equations with the flexibility of neural networks. By learning the drift and diffusion functions from observed market prices, our model inherently captures the volatility smile and other non-linearities.

\section{Methodology}
\subsection{Neural SDE Architecture}
We model the asset price $S_t$ as a stochastic process governed by a learned Itô SDE:

\begin{equation}
dS_t = \mu(t, S_t) dt + \sigma(t, S_t) dW_t
\end{equation}

Where:
\begin{itemize}
    \item $\mu(t, S_t)$: The drift function, approximated by a neural network.
    \item $\sigma(t, S_t)$: The diffusion (volatility) function, approximated by a neural network.
    \item $dW_t$: A Brownian motion increment.
\end{itemize}

\subsection{Training}
The model was implemented in PyTorch using the \texttt{torchsde} library. Key training details include:
\begin{align*}
    \text{Hardware} & : \text{NVIDIA RTX 4050 Laptop GPU} \\
    \text{Optimization} & : \text{Adjoint Sensitivity Method (O(1) Memory)} \\
    \text{Epochs} & : 50 \\
    \text{Data} & : \text{SPY, QQQ, IWM (Yahoo Finance)}
\end{align*}

\section{Experiments \& Results}
We benchmarked the trained Neural SDE against the standard Black-Scholes model using a hold-out test set of 347 real-world options (SPY).

\subsection{Quantitative Analysis}
The Neural SDE achieved significantly lower error metrics across all moneyness levels.

\begin{table}[h]
\centering
\caption{Comparison of Pricing Accuracy (347 Test Samples)}
\label{tab:results}
\begin{tabular}{lcc}
\toprule
\textbf{Metric} & \textbf{Neural SDE (Ours)} & \textbf{Black-Scholes (Baseline)} \\
\midrule
Mean Squared Error (MSE) & \textbf{0.26} & 389.87 \\
Mean Absolute Error (MAE) & \textbf{0.30} & 5.53 \\
Avg \$ Error & $\sim$\$0.50 & $\sim$\$19.70 \\
\bottomrule
\end{tabular}
\end{table}

\subsection{Qualitative Analysis}
\begin{itemize}
    \item \textbf{Volatility Smile}: The Neural SDE correctly priced deep Out-of-the-Money (OTM) and In-the-Money (ITM) options, whereas Black-Scholes systematically underpriced them due to its constant volatility assumption.
    \item \textbf{Adaptability}: The model successfully adapted to different time horizons (1 week to 1 year) without manual recalibration.
\end{itemize}

\section{Live Market Application}
The model was further applied to the Indian Stock Market (Nifty 50), demonstrating a real-time capability to generate trading signals based on the divergence between the model's computed fair value and the observed market price.

\section{Conclusion}
Neural SDEs represent a paradigm shift in derivatives pricing. By learning the stochastic dynamics from data, they outperform rigid analytical formulas by a substantial margin in mean squared error. Future work will extend this framework to exotic options and hedging strategies (Deep Hedging).

\end{document}
